\documentclass[11pt,a4paper]{article}
\usepackage[utf8]{inputenc}
\usepackage[T1]{fontenc}
\usepackage{amsmath,amssymb,amsthm,mathtools}
\usepackage[margin=2.5cm]{geometry}
\usepackage[hypertexnames=false]{hyperref}
\usepackage{booktabs}

\newtheorem{proposition}{Proposition}
\newtheorem{remark}{Remark}
\DeclareMathOperator*{\argmin}{arg\,min}
\DeclareMathOperator*{\argmax}{arg\,max}
\newcommand{\R}{\mathbb{R}}
\newcommand{\Hs}{\mathcal{H}}
\newcommand{\one}{\mathbf{1}}
\newcommand{\inner}[2]{\langle #1,\, #2 \rangle}

\title{\textbf{Project 29}\\[2mm]
\large Multi-class SVM, a Lagrangian dual of the representer identity,\\
and a bundle method}
\author{Tolgonai Nasipbek kyzy (626449) \and Hamse Hassan Adnan (683187)}
\date{}

\begin{document}
\maketitle

\begin{abstract}
\noindent
(M) is a kernel Support Vector Machine for multi-class classification built with
the standard One-vs-Rest construction. (A1) is a Lagrangian dual approach applied
to the \emph{dual} formulation of the SVM: we reintroduce the weight vector $w$ as
an explicit variable, so that the representer identity $w=\sum_i\alpha_iy_i\phi(x_i)$
becomes a constraint, and we dualise \emph{that} identity. The resulting Lagrangian
relaxation is a linear program over a box plus one equality constraint --- a
continuous knapsack, solved exactly by sorting in $O(m\log m)$ --- so the dual
function is finite everywhere, genuinely non-differentiable, and defined on a space
of dimension $n$ (linear kernel) or $m$ (general kernel). It is maximised by a
proximal bundle method, whose master problem is handed to an off-the-shelf QP solver.
(A2) is a general-purpose convex QP solver applied to the standard SVM dual.
\end{abstract}

\section{The standard SVM}
\label{sec:svm}

\subsection{Primal}
Let $\{(x_i,y_i)\}_{i=1}^m$ with $x_i\in\R^n$ and $y_i\in\{-1,+1\}$, and let
$\phi:\R^n\to\Hs$ be a feature map into a Hilbert space with reproducing kernel
$\kappa(x,x')=\inner{\phi(x)}{\phi(x')}$. Write $\phi_i:=\phi(x_i)$. The soft-margin
($C$-)SVM is
\begin{equation}
\label{eq:primal}
\tag{P}
\begin{aligned}
\min_{w\in\Hs,\,b\in\R,\,\xi\in\R^m}\quad & \tfrac12\|w\|^2 + C\sum_{i=1}^m \xi_i\\
\text{s.t.}\quad & y_i\bigl(\inner{w}{\phi_i}+b\bigr)\ \ge\ 1-\xi_i, \qquad i=1,\dots,m,\\
& \xi_i\ \ge\ 0, \qquad i=1,\dots,m .
\end{aligned}
\end{equation}
The first term maximises the geometric margin $1/\|w\|$ between the two supporting
hyperplanes $\inner{w}{\phi(x)}+b=\pm1$; the second penalises, linearly and with
weight $C>0$, the amount by which each sample fails to reach its own margin.
Problem~\eqref{eq:primal} is a convex QP; it is strictly convex in $w$ but only
convex in $(b,\xi)$.

Eliminating $\xi$ gives the equivalent unconstrained, non-smooth form
\begin{equation}
\label{eq:hinge}
\min_{w,b}\ \tfrac12\|w\|^2 + C\sum_{i=1}^m\bigl[1-y_i(\inner{w}{\phi_i}+b)\bigr]^+ ,
\qquad [t]^+:=\max\{0,t\},
\end{equation}
which we shall meet again, unexpectedly, in Section~\ref{sec:whatitis}.

\subsection{Wolfe dual}
Attach $\alpha_i\ge0$ to the margin constraints and $\mu_i\ge0$ to $\xi_i\ge0$:
\[
L(w,b,\xi;\alpha,\mu)=\tfrac12\|w\|^2+C\one^{\!\top}\xi
+\sum_{i=1}^m\alpha_i\bigl(1-\xi_i-y_i(\inner{w}{\phi_i}+b)\bigr)-\mu^{\!\top}\xi .
\]
Stationarity in $(w,b,\xi)$ gives the three conditions
\begin{equation}
\label{eq:representer}
w=\sum_{i=1}^m \alpha_i y_i \phi_i,
\qquad
\sum_{i=1}^m \alpha_i y_i = 0,
\qquad
\alpha_i+\mu_i=C .
\end{equation}
The first is the \emph{representer identity}; it is the hinge on which this whole
project turns --- Section~\ref{sec:choice} dualises precisely this condition.

Substituting the three conditions back into $L$ makes every term collapse. Writing
$K_{ij}=\kappa(x_i,x_j)$, $Y=\operatorname{diag}(y)$ and $Q=YKY$:
\begin{align*}
\tfrac12\|w\|^2 &= \tfrac12\sum_{i,j}\alpha_i\alpha_jy_iy_jK_{ij}
 = \tfrac12\alpha^{\!\top}Q\alpha, \\[1mm]
-\sum_i\alpha_iy_i\inner{w}{\phi_i}
 &= -\Bigl\langle w,\ \sum_i\alpha_iy_i\phi_i\Bigr\rangle = -\|w\|^2
 = -\,\alpha^{\!\top}Q\alpha, \\[1mm]
-\,b\sum_i\alpha_iy_i &= 0
 &&\text{(by the $b$ condition)},\\[1mm]
\sum_i\xi_i\bigl(C-\alpha_i-\mu_i\bigr) &= 0
 &&\text{(by the $\xi$ condition)},
\end{align*}
leaving $\one^{\!\top}\alpha-\tfrac12\alpha^{\!\top}Q\alpha$, to be maximised. Two
remarks on how the constraints of the dual arise, because both are easy to state
loosely:
\begin{itemize}
\item $\mu_i\ge0$ and $\mu_i=C-\alpha_i$ together give $\alpha_i\le C$, which with
      $\alpha_i\ge0$ is the box $0\le\alpha\le C\one$;
\item $b$ enters $L$ only through the term $-b\sum_i\alpha_iy_i$, and $b$ is
      \emph{free}, so $\inf_b L=-\infty$ unless $\sum_i\alpha_iy_i=0$. The
      $b$-stationarity condition therefore does not disappear when it is used to
      cancel that term: it survives as the constraint $y^{\!\top}\alpha=0$, which
      is exactly the statement that $\operatorname{dom}$ of the dual function
      excludes every $\alpha$ violating it.
\end{itemize}
Writing the result as a minimisation of the negative, the dual of~\eqref{eq:primal}
is the convex QP
\begin{equation}
\label{eq:dual}
\tag{D}
\min_{\alpha\in\R^m}\ \ \tfrac12\,\alpha^{\!\top} Q\,\alpha-\one^{\!\top}\alpha
\qquad\text{s.t.}\qquad \alpha\in D:=\bigl\{\alpha:\ y^{\!\top}\alpha=0,\ \ 0\le\alpha\le C\one\bigr\}.
\end{equation}
$D$ is a non-empty ($0\in D$) compact polytope and $Q\succeq0$, so~\eqref{eq:dual}
always has a solution. Slater's condition holds for~\eqref{eq:primal}, hence there
is no duality gap and the KKT conditions
\begin{equation}
\label{eq:kkt}
\alpha_i=0\ \Rightarrow\ y_if(x_i)\ge1,
\qquad
0<\alpha_i<C\ \Rightarrow\ y_if(x_i)=1,
\qquad
\alpha_i=C\ \Rightarrow\ y_if(x_i)\le1
\end{equation}
characterise optimality, where $f(x)=\sum_i\alpha_iy_i\kappa(x_i,x)+b$. Samples with
$\alpha_i>0$ are the \emph{support vectors}; those with $0<\alpha_i<C$ are the
\emph{free} support vectors, and $b$ is recovered from any of them via
$b=y_i-\sum_j\alpha_jy_j\kappa(x_j,x_i)$. The count of free support vectors will
turn out to control the cost of (A1); see Section~\ref{sec:cost}.

\subsection{Recovering the primal variables}
\label{sec:recovery}
Given an optimal $\alpha^\star$ of~\eqref{eq:dual}, the primal variables come back as
\[
w^\star=\sum_i\alpha^\star_iy_i\phi_i,
\qquad
\xi^\star_i=\bigl[\,1-y_i(\inner{w^\star}{\phi_i}+b^\star)\,\bigr]^+ ,
\]
the first directly from the representer identity and the second from primal
feasibility of the margin constraint together with complementarity on $\xi_i\ge0$
(note the factor $y_i$ inside: $\xi_i$ measures the violation of \emph{the sample's
own} margin, not of the hyperplane equation).

The bias $b^\star$ deserves more care, and it is where this report departs from the
usual treatment. The textbook recipe reads $b^\star$ off any free support vector via
\eqref{eq:kkt}, and averages over them for stability; but the average is only as
good as the set of indices identified as free, so a solver that is merely close to
optimal produces a $b$ that degrades the measured primal--dual gap. The better
observation --- made, for instance, in the interior-point treatment of the same
model in Project~25 --- is that $b$ enters the Lagrangian only through
$-b\sum_i\alpha_iy_i$, so \emph{$b$ is the multiplier of the equality constraint
$y^{\!\top}\alpha=0$}. Any algorithm that produces that multiplier explicitly can
use it directly, with no support-vector bookkeeping at all.

For an interior-point method this multiplier is one of the iterates. What is worth
noting here is that our relaxation supplies it too, and for free: the multiplier
$\theta$ of the single equality inside the continuous knapsack of
Section~\ref{sec:knapsack} \emph{is} $b$, as \eqref{eq:theta_is_b} proves by LP
duality. So (A1) returns the bias at every oracle call rather than only at
optimality, and the identity holds by construction rather than by an argument about
which $\alpha_i$ are free. We verify this numerically against a general-purpose QP
solver to $10^{-7}$.

\subsection{Kernels}
Only inner products appear in~\eqref{eq:dual}, so $\phi$ need never be formed. We use
the linear kernel $\kappa(x,x')=x^{\!\top}x'$, the Gaussian kernel
$\kappa(x,x')=\exp(-\gamma\|x-x'\|^2)$, and the polynomial kernel
$\kappa(x,x')=(\gamma\,x^{\!\top}x'+c_0)^d$.

\section{One-vs-Rest}
\label{sec:ovr}

With $q$ classes and labels $c_i\in\{1,\dots,q\}$, One-vs-Rest (OvR) solves $q$
independent binary problems. For class $k$ set
\[
y^k_i=\begin{cases}+1,&c_i=k\\ -1,&c_i\neq k\end{cases}
\]
and solve~\eqref{eq:primal}/\eqref{eq:dual} with $y=y^k$, obtaining
$f^k(x)=\sum_i\alpha^k_iy^k_i\kappa(x_i,x)+b^k$. A new point is assigned to
$\argmax_k f^k(x)$. Note that:

\begin{enumerate}
\item The $q$ problems share \emph{the same} kernel matrix $K$; they differ only
      through the label vector, since $Q^k=Y^kKY^k$ and $D^k$ depends on $k$ only
      through the single equality $ (y^k)^{\!\top}\alpha=0$. $Q^k$ is never formed.
\item The classifiers are trained independently but compared at prediction time on
      a common scale that nothing in the training enforces; this produces the
      familiar OvR pathologies, ties ($f^k(x)>0$ for two $k$) and rejections
      ($f^k(x)<0$ for all $k$).
\item Each binary problem is unbalanced, roughly $1:(q-1)$.
\end{enumerate}
Point~1 is the basis of the discussion in Section~\ref{sec:transfer}, which the
project statement explicitly requires.

\section{Which constraints to dualise}
\label{sec:choice}

\subsection{Three requirements}
For a Lagrangian dual approach to be worth building, the relaxation must satisfy:
\begin{description}
\item[(R1)] the relaxed problem must be \emph{substantially easier} than the
  original --- ideally closed-form or combinatorial;
\item[(R2)] the dual must be \emph{multi-dimensional}, otherwise a one-dimensional
  search solves it and bundle machinery is pointless;
\item[(R3)] the dual should be \emph{genuinely non-differentiable}, otherwise a
  smooth method is preferable.
\end{description}
These are exactly the three tests our earlier proposals failed, and they are worth
stating because they make the correct choice inevitable.

Relaxing the single equality $y^{\!\top}\alpha=0$ in~\eqref{eq:dual} violates (R1)
and (R2): what is left is
\[
\min\Bigl\{\tfrac12\alpha^{\!\top}Q\alpha-\one^{\!\top}\alpha+\theta\,y^{\!\top}\alpha
\ :\ 0\le\alpha\le C\one\Bigr\},
\]
a box-constrained QP with the \emph{same dense Hessian} --- the unbiased SVM, no
easier than the original --- and the dual is univariate. Relaxing the box
$\alpha\le C\one$ fares no better: the relaxation is again a QP in $Q$.

The reason is structural, and it is the point we had missed: in~\eqref{eq:dual}
\emph{no constraint is the difficulty}. The difficulty is $Q$ itself, a dense
$m\times m$ matrix. So the object to dualise is not a constraint of~\eqref{eq:dual}
at all --- it is the identity that \emph{created} $Q$.

\subsection{Dualising the representer identity}
$Q$ appears in~\eqref{eq:dual} only because $w$ was eliminated through the
representer identity~\eqref{eq:representer}. Put $w$ back:
\begin{equation}
\label{eq:split}
\tag{D$'$}
\min_{\alpha,\,w}\ \ \tfrac12\|w\|^2-\one^{\!\top}\alpha
\qquad\text{s.t.}\qquad
w=\sum_{i=1}^m\alpha_iy_i\phi_i,
\qquad \alpha\in D .
\end{equation}
Substituting the constraint recovers~\eqref{eq:dual} exactly, so \eqref{eq:split}
and \eqref{eq:dual} are the same problem. Now dualise \emph{only} the linking
identity, with a free multiplier $\lambda\in\Hs$:
\[
L(\alpha,w;\lambda)=\tfrac12\|w\|^2-\one^{\!\top}\alpha
+\Bigl\langle \lambda,\ w-\sum_i\alpha_iy_i\phi_i \Bigr\rangle .
\]
The two groups of variables separate. Minimising over $w\in\Hs$ is unconstrained
and gives $w(\lambda)=-\lambda$ with value $-\tfrac12\|\lambda\|^2$; what remains is
linear in $\alpha$ over $D$. Hence
\begin{equation}
\label{eq:phi}
\boxed{\ \
\psi(\lambda)\ =\ -\tfrac12\|\lambda\|^2
\ +\ \min_{\alpha\in D}\ c(\lambda)^{\!\top}\alpha,
\qquad
c_i(\lambda)=-\bigl(1+y_i\inner{\lambda}{\phi_i}\bigr),
\ \ }
\end{equation}
and the Lagrangian dual of~\eqref{eq:split} is $\max_{\lambda\in\Hs}\psi(\lambda)$.

In the vocabulary of the course's treatment of partial Lagrangian relaxation, the
representer identity is the \emph{complicating} constraint --- it is what couples
$w$ to $\alpha$ and so what creates the dense $Q$ --- while $\alpha\in D$ are the
\emph{easy} constraints, kept in the relaxed problem because they are exactly what
makes it a knapsack. The usual warning that $\psi(\lambda)$ may be $-\infty$ for
some $\lambda$, and that one then has to neutralise the offending directions of zero
curvature, does not apply here: $D$ is compact, so the relaxed problem always has a
finite optimum and $\psi$ is finite on all of $\Hs$.

\subsection{The relaxed problem is a continuous knapsack}
\label{sec:knapsack}
The inner problem in~\eqref{eq:phi} is
\begin{equation}
\label{eq:kp}
\tag{KP}
\min\ c^{\!\top}\alpha \qquad\text{s.t.}\qquad y^{\!\top}\alpha=0,\quad 0\le\alpha\le C\one,
\end{equation}
a linear program over a box with a \emph{single} equality constraint: the classical
\emph{continuous knapsack} problem. Relaxing that one equality with a scalar
$\theta$ makes the box minimiser explicit,
\[
\alpha_i(\theta)=\begin{cases}
C, & c_i-\theta y_i<0,\\
0, & c_i-\theta y_i>0,\\
\text{any value in }[0,C], & c_i-\theta y_i=0 ,
\end{cases}
\]
and $h(\theta):=y^{\!\top}\alpha(\theta)$ is non-decreasing, so a $\theta^\star$ with
$h(\theta^\star)=0$ exists and is found by sorting the $m$ breakpoints $y_ic_i$ and
scanning (or by median search in $O(m)$ expected time). Splitting the tied
components to make the equality hold exactly produces an $\alpha$ that minimises the
Lagrangian of~\eqref{eq:kp} \emph{and} is feasible, hence optimal. Forming $c$ costs
one kernel product; the rest costs $O(m\log m)$. This settles (R1): a dense QP has
been replaced by a sort.

\subsection{Properties of $\psi$}
\begin{proposition}[finiteness]
$D$ is a non-empty compact polytope, so the minimum in~\eqref{eq:phi} is attained
and finite for every $\lambda$. Hence $\psi:\Hs\to\R$ is finite everywhere,
concave and upper semicontinuous.
\end{proposition}
This is unusually convenient: no feasibility (``infinite'') cuts are ever needed,
and every oracle call returns a valid lower bound on the optimal value.

\begin{proposition}[structure and non-differentiability]
$\psi=\psi_{\mathrm{sm}}+\psi_{\mathrm{pl}}$ where
$\psi_{\mathrm{sm}}(\lambda)=-\tfrac12\|\lambda\|^2$ is smooth and known in
closed form, and $\psi_{\mathrm{pl}}(\lambda)=\min_{\alpha\in D}c(\lambda)^{\!\top}\alpha$
is a minimum of finitely many affine functions of $\lambda$ (one per vertex of $D$),
hence concave polyhedral. $\psi$ fails to be differentiable exactly where
\eqref{eq:kp} has multiple optima, i.e. where some reduced cost vanishes,
$c_i-\theta^\star y_i=0$.
\end{proposition}
Unwinding that condition with $w=-\lambda$ gives $y_i\inner{w}{\phi_i}+\theta^{\star}y_i=1$:
the kinks of $\psi$ sit \emph{exactly} on the support-vector condition
$y_if(x_i)=1$ of~\eqref{eq:kkt}. Since any non-degenerate SVM has free support
vectors, the maximiser of $\psi$ \emph{is} a point of non-differentiability. This
settles (R3): the non-smoothness is not incidental, it is the solution.

\begin{proposition}[no duality gap]
\eqref{eq:split} is convex, its coupling constraint is affine and $D$ is compact;
therefore $\max_\lambda\psi(\lambda)=v(\mathrm{D}')=v(\mathrm{D})$, and the
maximiser is $\lambda^\star=-w^\star$.
\end{proposition}
The multiplier of the relaxed identity is thus, up to sign, the primal weight
vector itself. A supergradient of $\psi$ at $\lambda$ is the violation of the
relaxed constraint,
\begin{equation}
\label{eq:supergrad}
g(\lambda)=w(\lambda)-\sum_i\alpha_i(\lambda)y_i\phi_i=-\lambda-\sum_i\alpha_i(\lambda)y_i\phi_i ,
\end{equation}
available at no extra cost from the oracle. (R2) holds as well: $\lambda$ ranges over
$\Hs$, of dimension $n$ for the linear kernel and $m$ after the reduction of
Section~\ref{sec:kernelised} in general.

\subsection{What this dual actually is}
\label{sec:whatitis}
It is worth identifying $\max_\lambda\psi$ in closed form. Substituting
$w=-\lambda$ in~\eqref{eq:phi},
\[
\max_{\lambda}\psi(\lambda)
\iff
-\min_{w}\Bigl\{\tfrac12\|w\|^2+\max_{\alpha\in D}\sum_i\alpha_i\bigl(1-y_i\inner{w}{\phi_i}\bigr)\Bigr\},
\]
and dualising the single equality inside the inner maximisation with a scalar $\theta$
gives, by LP duality,
\begin{equation}
\label{eq:theta_is_b}
\max_{\alpha\in D}\sum_i\alpha_i\bigl(1-y_ig_i\bigr)
=\min_{\theta}\ C\sum_i\bigl[1-y_i(g_i+\theta)\bigr]^+ ,
\qquad g_i:=\inner{w}{\phi_i} .
\end{equation}
So the Lagrangian dual of~\eqref{eq:split} is the hinge-loss primal~\eqref{eq:hinge},
and \emph{the multiplier $\theta$ of the knapsack equality is the SVM bias $b$}
(up to the sign convention of Section~\ref{sec:knapsack}) --- so $b$ falls out of the
oracle for free, with no need for the usual ``average over the free support vectors''
recipe. That the dual of the dual is the primal is of course biconjugacy; what is
useful here is the algorithmic consequence, namely that (A1) is a bundle method on a
function for which published convergence results exist: the cutting-plane and bundle
methods for regularised risk minimisation of Joachims (KDD~2006), Teo et al.
(JMLR~2010) and Franc \& Sonnenburg (ICML~2008) converge in $O(1/\varepsilon)$
iterations, improving to linear under the strong convexity of $\tfrac12\|w\|^2$.

\subsection{Relation to the nonsmooth formulation of the course slides}
\label{sec:bmrm}
The SVM appears in the course's own deck on nondifferentiable unconstrained
optimization as the motivating example, in the form
$\min_{w,b}\ \|w\|^2+C\sum_{i\in I}\max\{1-y_i(\inner{w}{x_i}-b),0\}$, whose
nondifferentiability is precisely what motivates bundle methods. Since
Section~\ref{sec:whatitis} shows that $\max_\lambda\psi$ is that same function, we
must say clearly what dualising the representer identity adds over simply running a
bundle method on it --- which is also what the published cutting-plane methods for
regularised risk minimisation do (Joachims KDD~2006; Teo et al.\ JMLR~2010; Franc \&
Sonnenburg ICML~2008).

The answer is \emph{not} that the bias makes a direct approach impossible: $f(w,b)$
is convex in $(w,b)$ jointly and a subgradient in both variables is immediate, so
one can perfectly well put $b$ in the bundle. The answer is that our relaxation
minimises over $b$ \emph{exactly}, at every oracle call, instead of approximating
the $b$ direction with cutting planes. The knapsack multiplier $\theta$
of~\eqref{eq:theta_is_b} \emph{is} that exact partial minimisation, and the sort of
Section~\ref{sec:knapsack} is how it is computed. Equivalently: we bundle
$g(w)=\min_b f(w,b)$ rather than $f(w,b)$, and $g$ is the better-conditioned object
because one of its directions has been eliminated in closed form.

We measured the difference against the direct baseline, keeping everything else
identical --- the quadratic term treated exactly, a cutting-plane model of the rest,
the same proximal term and the same master solver --- so that the only difference is
whether $b$ is a bundle variable:
\begin{center}
\begin{tabular}{llrrr}
\toprule
dataset & kernel & $b$ in the bundle & $b$ eliminated (ours) & ratio\\
\midrule
\textsc{wine}   & linear   & 137 & 29 & $4.7\times$\\
\textsc{wine}   & Gaussian & \multicolumn{1}{c}{400 (gap $6\cdot10^{-5}$)} & 133 & did not converge\\
\textsc{iris}   & Gaussian & 96 & 24 & $4.0\times$\\
\textsc{cancer} & linear   & 220 & 52 & $4.2\times$\\
\bottomrule
\end{tabular}
\end{center}
(iterations to a relative gap of $10^{-8}$; the Gaussian \textsc{wine} baseline was
still at $6\cdot10^{-5}$ after the 400-iteration cap, where ours reached
$1.1\cdot10^{-8}$ in 133.) A consistent factor of four to five, and one instance
where the direct method does not reach tolerance at all.

That is the concrete content of the dualisation, and we claim exactly this much: the
family of cuts and the function being optimised are the same as in the direct
approach, and what relaxing the representer identity adds is (i) the exact
elimination of $b$ at every oracle call, worth the factor above, (ii) the SVM dual
variables, through the aggregated solution of Section~\ref{sec:dw}, and (iii) a
two-sided bound on $v(\mathrm{D})$ rather than a bound on one's own objective only.

\section{(A1): the bundle method}
\label{sec:bundle}

\subsection{Model and master problem}
Because $-\tfrac12\|\lambda\|^2$ is known in closed form, we linearise only the
polyhedral part --- a \emph{generalized} (structured) bundle method. Each oracle call
returns some $\alpha^j\in D$, which gives the affine upper model
\[
\psi_{\mathrm{pl}}(\lambda)\ \le\ c(\lambda)^{\!\top}\alpha^j
= -\,\one^{\!\top}\alpha^j-\inner{\lambda}{\textstyle\sum_i\alpha^j_iy_i\phi_i}
=: -a_j-\inner{\lambda}{\Phi^{\!\top}Y\alpha^j} .
\]
With a stability centre $\bar\lambda$ and proximal parameter $t>0$ the master is
\begin{equation}
\label{eq:master}
\max_{\lambda\in\Hs}\ \
-\tfrac12\|\lambda\|^2
+\min_{j=1,\dots,J}\bigl\{-a_j-\inner{\lambda}{b_j}\bigr\}
-\tfrac{1}{2t}\|\lambda-\bar\lambda\|^2 ,
\qquad b_j:=\Phi^{\!\top}Y\alpha^j .
\end{equation}
Exchanging $\min$ and $\max$ (Sion) and eliminating $\lambda$ gives its dual, a QP
over the unit simplex $\Delta^J$:
\begin{equation}
\label{eq:masterdual}
\min_{\vartheta\in\Delta^J}\ \
\frac{1}{2\rho}\Bigl\|\tfrac{\bar\lambda}{t}-\textstyle\sum_j\vartheta_jb_j\Bigr\|^2
-\ \textstyle\sum_j\vartheta_j a_j ,
\qquad \rho:=1+\tfrac1t,
\end{equation}
after which $\lambda^{\mathrm{new}}=\bigl(\bar\lambda/t-\sum_j\vartheta_jb_j\bigr)/\rho$.
Only the $J\times J$ Gram matrix $G_{jl}=\inner{b_j}{b_l}$ and the vector
$p_j=\inner{b_j}{\bar\lambda}$ are needed; both are maintained incrementally, so a
master costs $O(J^2)$ plus one kernel product per new cut. We solve it with an
off-the-shelf QP solver, which the project statement allows.

\subsection{Primal recovery, and the Dantzig--Wolfe identity}
\label{sec:dw}
Let $\hat\alpha:=\sum_j\vartheta_j\alpha^j$ be the \emph{aggregated solution} built from the master problem's dual multipliers. Since $\vartheta\in\Delta^J$ and
$D$ is convex, $\hat\alpha\in D$ --- it is a \emph{feasible} solution
of~\eqref{eq:dual} at every iteration. Moreover $\sum_j\vartheta_jb_j=\Phi^{\!\top}Y\hat\alpha$
and $\sum_j\vartheta_ja_j=\one^{\!\top}\hat\alpha$, so for $t\to\infty$
\eqref{eq:masterdual} becomes
\[
\min_{\vartheta\in\Delta^J}\ \tfrac12\,\hat\alpha^{\!\top}Q\hat\alpha-\one^{\!\top}\hat\alpha ,
\]
i.e. \emph{the SVM dual \eqref{eq:dual} restricted to $\mathrm{conv}\{\alpha^1,\dots,\alpha^J\}\subseteq D$}.
The method is therefore exactly Dantzig--Wolfe decomposition of~\eqref{eq:dual} with
\eqref{eq:kp} as pricing problem, and bundle, cutting planes on~\eqref{eq:hinge} and
column generation on~\eqref{eq:dual} are three views of one algorithm. Two practical
consequences:
\begin{itemize}
\item the master returns the primal solution $\hat\alpha\to\alpha^\star$, so nothing
      extra is needed to recover the classifier;
\item $\psi(\lambda)\le v(\mathrm{D})\le \tfrac12\hat\alpha^{\!\top}Q\hat\alpha-\one^{\!\top}\hat\alpha$
      at \emph{every} iteration, a \textbf{certified gap} --- which decomposition
      algorithms for the SVM such as SMO do not provide.
\end{itemize}

\subsection{Kernelised form}
\label{sec:kernelised}
For a general kernel $\Hs$ may be infinite-dimensional, so write
$\lambda=\sum_j\beta_j\phi_j$. This is without loss of generality: a component of
$\lambda$ orthogonal to $\mathrm{span}\{\phi_i\}$ leaves $c(\lambda)$ unchanged and
only increases $\|\lambda\|^2$, so it is zero at the optimum. Then
$\inner{\lambda}{\phi_i}=(K\beta)_i$ and $\|\lambda\|^2=\beta^{\!\top}K\beta$, and one
oracle call costs one matrix--vector product with $K$ plus a sort. Note that the
correct stabilising term is the \emph{RKHS} norm
$\|\lambda-\bar\lambda\|^2=(\beta-\bar\beta)^{\!\top}K(\beta-\bar\beta)$ rather than
the Euclidean $\|\beta-\bar\beta\|^2$; it is the $K$-metric that makes
\eqref{eq:masterdual} come out in the form above.

\subsection{Cost, and how many cuts are needed}
\label{sec:cost}
Per iteration: one oracle call ($O(m^2)$ for the kernel product, $O(m\log m)$ for the
sort) plus one master ($O(J^2)$ in the QP, $O(Jm)$ to update $G$). The interesting
question is how large $J$ must grow. The oracle returns essentially \emph{vertices}
of $D$, that is $0/C$ vectors, and $\hat\alpha$ is a convex combination of them: a
component of $\alpha^\star$ sitting at a bound is produced by a single vertex, but
every \emph{free} support vector $0<\alpha^\star_i<C$ must be produced by mixing
vertices that disagree on $i$. We therefore expect, and measure, that
\begin{center}
\emph{the number of cuts the method needs is governed by the number of free support
vectors,}
\end{center}
and that truncating the bundle below that number makes the gap stall --- the classical
Dantzig--Wolfe tailing-off. This is confirmed sharply in our prototype: on
\textsc{wine}, the linear kernel yields $12$ free support vectors and the method
converges to a gap of $2\cdot10^{-8}$ in $26$ iterations, whereas the Gaussian kernel
($\gamma=0.2$) yields $91$ free support vectors out of $m=178$, and with a bundle
capped at $60$ cuts the gap stalls at $2\cdot10^{-4}$, while raising the cap to $200$
restores convergence in $153$ iterations. The practical rule is to size the bundle
above the expected number of free support vectors, which large $C$ and simple kernels
make small.

\section{(A2): a general-purpose solver}
For comparison, \eqref{eq:dual} is handed as a dense convex QP to an off-the-shelf
interior-point solver. This is the reference both for solution quality (we verify
that (A1) and (A2) return the same $\alpha^\star$, $b$ and test accuracy) and for
running time.

\section{Does solving one One-vs-Rest problem help the others?}
\label{sec:transfer}
The project statement asks this explicitly. The $q$ problems differ only through
$y^k$, so there are three candidate mechanisms, and they are not equally useful.

\paragraph{1. The kernel matrix (real).} $K$ is computed once and shared; $Q^k=Y^kKY^k$
is never formed, since the oracle only ever needs $K\beta$. Better still, the $q$
oracles can be run in lockstep so that their $q$ matrix--vector products become one
BLAS-3 product $K\,[\beta^1\cdots\beta^q]$. This is a genuine, structural saving of a
factor $q$ in the dominant cost.

\paragraph{2. Warm-starting the multiplier (measured, and it does not help).} Since
$\lambda^\star=-w^\star$, warm-starting class $k{+}1$ at class $k$'s optimum amounts
to assuming the two separating hyperplanes are similar. They are not: the positive
class changes completely between the two problems. In our experiments this is neutral
at best and can be clearly harmful (on \textsc{wine} with a Gaussian kernel it raised
the total oracle-call count by $72\%$), because the bundle inherits a stability centre
far from the new optimum and spends iterations undoing it.

\paragraph{3. Transferring cuts (measured, and it is roughly neutral).} Any point of
$D^l$ generates a valid cut for class $l$, and the $D^k$ differ only through the
equality. Writing $P_k$, $P_l$ for the two positive blocks and $R$ for the rest,
$\alpha\in D^k$ satisfies $\alpha(P_k)=\alpha(P_l)+\alpha(R)$, so
$(y^l)^{\!\top}\alpha=-2\,\alpha(R)$: the imbalance is known in closed form and can be
repaired in $O(m)$ by moving mass onto $P_l$, \emph{preserving the support pattern} of
the transferred cut. This costs no oracle call and no kernel product. Measured effect:
between $-4\%$ and $+2\%$ oracle calls --- within noise.

\medskip\noindent
The honest conclusion, which we will state as such in the report, is that for this
decomposition the transferable quantity is the \emph{kernel matrix}, not the dual
solution: mechanism~1 is worth a factor $q$, while mechanisms~2 and~3 are not worth
their complexity. That is a more useful answer than a plausible-sounding claim about
warm starts, and it is the one the measurements support.

\appendix
\section{Appendix: the shared-slack-budget model, and why we set it aside}
\label{app:budget}
For completeness we record the model proposed in our previous submission, now
derived properly, together with a defect we found in it.

\paragraph{The model.} Couple the $q$ One-vs-Rest problems by giving each sample a
single slack budget $B_i$ to be shared by all $q$ binary classifiers:
\begin{equation}
\label{eq:budget}
\tag{M-B}
\begin{aligned}
\min_{w,b,\xi}\quad & \sum_{k=1}^q\tfrac12\|w^k\|^2+C\sum_{i,k}\xi^k_i\\
\text{s.t.}\quad & y^k_i\bigl(\inner{w^k}{\phi_i}+b^k\bigr)\ \ge\ 1-\xi^k_i,
\qquad \xi^k_i\ \ge\ 0,
\qquad \sum_{k=1}^q \xi^k_i\ \le\ B_i .
\end{aligned}
\end{equation}

\paragraph{Its dual.} Dualising only the $mq$ margin constraints with
$\alpha^k_i\ge0$ separates the variables. Minimising over $w^k$ gives
$w^k=\sum_i\alpha^k_iy^k_i\phi_i$ and $-\tfrac12(\alpha^k)^{\!\top}Q^k\alpha^k$;
minimising over $b^k$ forces $\sum_i\alpha^k_iy^k_i=0$; and minimising over $\xi$
splits across $i$ into the tiny LP
$\min\{\sum_k(C-\alpha^k_i)\xi^k_i:\ \xi_i\ge0,\ \sum_k\xi^k_i\le B_i\}$,
solved by inspection: put the whole budget on the most negative coefficient, or use
none of it. Hence
\begin{equation}
\label{eq:budgetdual}
\psi(\alpha)=\sum_{i,k}\alpha^k_i-\tfrac12\sum_{k}(\alpha^k)^{\!\top}Q^k\alpha^k
-\sum_{i}B_i\Bigl[\max_k\bigl(\alpha^k_i-C\bigr)\Bigr]^+ ,
\end{equation}
on $\{\alpha\ge0:\ \sum_i\alpha^k_iy^k_i=0\ \ \forall k\}$. The budget replaces the
box $\alpha\le C$ of~\eqref{eq:dual} by a penalty, and the dual is non-differentiable
where two or more classes tie for a contested budget, or at the kink
$\max_k\alpha^k_i=C$ --- which is where the One-vs-Rest coupling actually happens.

\paragraph{The defect.} With a \emph{hard} budget, \eqref{eq:budget} can be
infeasible, and then its Lagrangian dual is unbounded. This is not a corner case: a
sample lying where two classifiers both score near zero needs $\xi^k_i\approx1$ for
each of them, so $B_i<2$ is typically already infeasible. On a $3$-class Gaussian
blob instance we measured the feasibility threshold at exactly $B_{\min}=2.0$, with
\eqref{eq:budgetdual} diverging for every $B<B_{\min}$ and strong duality holding to
$10^{-8}$ for every $B\ge B_{\min}$. A single outlier destroys the training problem.

\paragraph{The repair.} Soften the budget: $\sum_k\xi^k_i\le B_i+z_i$ with $z_i\ge0$
penalised by $Dz_i$. The $\xi$ subproblem stays an LP solved by inspection, the model
is always feasible, and boundedness of the $z$ term restores a box,
$0\le\alpha\le (C+D)\one$, making $\mathrm{dom}\,\psi$ compact. We verified
strong duality for the repaired model to $10^{-8}$ across a range of $B$, $D$, $q$
and $m$.

\paragraph{Why we nevertheless prefer the approach of Section~\ref{sec:choice}.}
Even repaired, \eqref{eq:budgetdual} retains the quadratic forms $Q^k$ \emph{inside}
the dual function, so the bundle's piecewise-linear model has to approximate a curved
function and the method degenerates into solving a QP by cutting planes. The
relaxation of Section~\ref{sec:choice}, by contrast, leaves an LP and keeps its only
curved term in closed form in the master. And \eqref{eq:budget} changes the model (M),
so any comparison against One-vs-Rest would confound a modelling change with an
algorithmic one, whereas Section~\ref{sec:choice} leaves (M) exactly the standard
One-vs-Rest SVM and changes only how it is solved.

\end{document}
